% task3_en.tex -- Task 3: Baseline Methods (included by main_en.tex; do not compile standalone)

\section*{Task 3: Baseline Methods}
\begin{itemize}
    \item \textbf{MLP}: Multi-layer perceptron. It uses no hypergraph structural information at all and only applies layer-wise non-linear transformations to node features, serving as the most basic baseline for measuring whether structural information yields any gain.
    \item \textbf{CEGCN}\cite{chien2022you,zhou2006learning,kipf2017semi}: A clique-expansion-based method. Each hyperedge is first converted into a clique over its incident nodes, turning the hypergraph into a pairwise graph, on which GCN is then applied to learn node representations.
    \item \textbf{CEGAT}\cite{chien2022you,zhou2006learning,velickovic2018graph}: Also converts the hypergraph into a pairwise graph via clique expansion, but employs GAT (graph attention network) to model the importance weights among nodes within a clique, so that the contributions of different neighbours are distinguished.
    
    \vspace{0.5em} % dashed grouping
    
    \item \textbf{HGNN}\cite{feng2019hypergraph}: Extends spectral convolution to hypergraphs. It exploits the hypergraph Laplacian and approximates the spectral filter with truncated Chebyshev polynomials, thereby capturing high-order correlations in complex data.
    \item \textbf{HyperGCN}\cite{yadati2019hypergcn}: Approximates each hyperedge with pairwise edges connecting the vertices inside it, reducing hypergraph learning to graph learning on the approximated graph.
    \item \textbf{HCHA}\cite{bai2021hypergraph}: Proposes end-to-end trainable hypergraph convolution and hypergraph attention operators. The convolution propagates information using high-order relations and the local clustering structure, while the attention dynamically adjusts hyperedge connections to aggregate task-relevant information.
    \item \textbf{LEGCN}\cite{yang2022semi}: Based on line expansion, it treats vertex--hyperedge pairs as nodes of a simple graph, thereby bijectively converting a hypergraph into a simple graph. Existing graph learning algorithms can thus be applied directly to hypergraphs without losing the symmetric co-occurrence structure.
    \item \textbf{HyperND}\cite{prokopchik2022nonlinear}: Constructs a non-linear diffusion process on the hypergraph that propagates features and labels along the hypergraph structure. The diffusion involves several non-linearities to strengthen the modelling capacity; its limit points serve as node embeddings, which are fed to a linear model for prediction.
    \item \textbf{PhenomNN}\cite{wang2023from}: Based on a family of expressive parametrized hypergraph-regularized energy functions, it defines node embeddings as the minimizers of these energies and jointly trains them with a parametrized classifier through supervised bilevel optimization.
    \item \textbf{SheafHyperGNN}\cite{duta2023sheaf}: Introduces the cellular sheaf framework to hypergraphs, modelling complex dynamics while preserving high-order connectivity. It generalizes the commonly used hypergraph Laplacian to a sheaf-structured version and constructs two neural networks that can infer and process sheaf structures on hypergraphs.
    \item \textbf{HJRL}\cite{yan2024hypergraph}: Proposes cross expansion, which converts both hypervertices and hyperedges into vertices of a standard graph and jointly embeds them in a shared representation space, while a hypergraph reconstruction objective preserves structural information.
    \item \textbf{DPHGNN}\cite{saxena2024dphgnn}: A hybrid framework for resource-constrained hypergraph scenarios. It introduces equivariant operator learning to capture low-order semantics, adopts a static update layer to provide spectral bias and relational features, and a dynamic update layer that fuses explicitly aggregated features into hypergraph message passing.
    \item \textbf{TF-HNN}\cite{tang2025training}: The first model that decouples hypergraph structure processing from model training, proposing a training-free message passing module, TF-MP-Module. The module removes learnable parameters and non-linear activations and compresses multi-layer propagation into a single step, thereby improving training efficiency.
    
    \vspace{0.5em} % dashed grouping
    
    \item \textbf{HNHN}\cite{dong2020hnhn}: A hypergraph convolutional network that applies non-linear activations on both hypernodes and hyperedges, together with a normalization scheme, so that the importance of high-order hyperedges and high-degree vertices can be adjusted flexibly per dataset.
    \item \textbf{UniGNN}\cite{huang2021unignn}: A unified message passing framework that generalizes standard GNNs to hypergraphs. It first aggregates the incident vertices with a permutation-invariant function to obtain hyperedge representations, and then aggregates the incident hyperedge representations to update each vertex.
    \item \textbf{AllSetTransformer}\cite{chien2022you,lee2019set}: Implements each layer as a composition of two multiset functions and incorporates the Set Transformer, which yields greater modelling flexibility and expressive power.
    \item \textbf{ED-HNN}\cite{wang2023equivariant}: Designs an architecture that approximates arbitrary continuous, permutation-equivariant hypergraph diffusion operators. It is realized by star expansion (a bipartite representation) combined with standard message passing neural networks, and supports scalable training with weights shared across layers.
    \item \textbf{HyperGT}\cite{liu2024hypergraph}: A Transformer-based HNN for capturing global correlations between nodes and hyperedges. It introduces incidence-matrix-based positional encoding and a structural regularization term to preserve local structure, jointly modelling global interactions and local connectivity patterns.
    
    \vspace{0.5em} % dashed grouping
    
    \item \textbf{EHNN}\cite{kim2022equivariant}: The first framework that realizes general hypergraph-equivariant GNNs. It establishes a link between sparse hypergraphs and dense fixed-order tensors, designs maximally expressive equivariant linear layers, and guarantees scalability and generalization to arbitrary hyperedge orders through hypernetwork parameter sharing.
    \item \textbf{T-HyperGNN}\cite{wang2024thypergnns,pena2023thgsp}: Integrates tensor hypergraph signal processing (t-HGSP) and encodes the hypergraph structure with tensors. Node interactions are modelled by multiplicative interaction tensors, lifting aggregation from linear operations to high-order polynomial mappings, and tensor sparsity is exploited for efficient tensor message passing.
\end{itemize}
